\begin{trapbox}
	\begin{description}[style=nextline, leftmargin=2.8em, labelsep=0.5em, itemsep=0.35em]
		\item[\textbf{\textcolor{CTrap}{易错点一（忽略新元范围）}}] 在换元后，未明确新变量的取值范围，导致值域扩大或缩小。
		\item[\textbf{\textcolor{CTrap}{正确理解（同步确定范围）}}] 每一步换元都必须根据原变量和新变量的对应关系，同步写出新变量的范围。例如 $t=\sqrt{ax+b}$ 必须注明 $t\ge0$。
		\item[\textbf{\textcolor{CTrap}{错因分析（忘记限制条件）}}] 学生常直接代入表达式求值域，却忽略根式本身的非负性或三角函数的定义区间，造成值域错误。
	\end{description}

	\medskip
	\begin{description}[style=nextline, leftmargin=2.8em, labelsep=0.5em, itemsep=0.35em]
		\item[\textbf{\textcolor{CTrap}{易错点二（三角换元的角域选取）}}] 对于 $\sqrt{a^2-x^2}$ 采用 $x=a\sin\theta$ 时，若选取 $\theta$ 范围不当，可能导致 $\cos\theta$ 符号不定，根式开方出现正负混淆。
		\item[\textbf{\textcolor{CTrap}{正确理解（为保证非负选取合适角域）}}] 规定 $\theta\in[-\frac{\pi}{2},\frac{\pi}{2}]$，此时 $\cos\theta\ge0$，保证 $\sqrt{a^2-x^2}=a\cos\theta$ 成立。
		\item[\textbf{\textcolor{CTrap}{错因分析（未固定符号盲目开方）}}] 学生不熟悉恒等变形后的开方规则，直接写 $\sqrt{1-\sin^2\theta}=\cos\theta$ 而未考虑符号，导致结果正负号错误。
	\end{description}

	\medskip
	\begin{description}[style=nextline, leftmargin=2.8em, labelsep=0.5em, itemsep=0.35em]
		\item[\textbf{\textcolor{CTrap}{易错点三（正切与正割换元混淆）}}] 对于 $\sqrt{x^2+a^2}$ 与 $\sqrt{x^2-a^2}$ 两种形式，学生容易混淆使用的换元方式。
		\item[\textbf{\textcolor{CTrap}{正确理解（区分平方和与平方差）}}] $\sqrt{x^2+a^2}$ 使用 $x=a\tan\theta$；$\sqrt{x^2-a^2}$ 使用 $x=a\sec\theta$。记忆方法：和用正切，差用正割。
		\item[\textbf{\textcolor{CTrap}{错因分析（没有抓住恒等式特征）}}] 关键在于识别根号内是“平方相加”还是“平方相减”。相加用 $1+\tan^2\theta=\sec^2\theta$，相减用 $\sec^2\theta-1=\tan^2\theta$，混淆则导致推理错误。
	\end{description}
\end{trapbox}
